% =====================================================================
% BMED365: Neurosymbolic AI and Agentic AI
% Beamer-presentasjon - Topic List S01-S11
% =====================================================================
\documentclass[aspectratio=169, 10pt]{beamer}

% =====================================================================
% PAKKER
% =====================================================================
\usepackage[utf8]{inputenc}
\usepackage[T1]{fontenc}
\usepackage[norsk]{babel}
\usepackage{graphicx}
\usepackage{tikz}
\usetikzlibrary{shapes.geometric, arrows, positioning, calc}
\usepackage{booktabs}
\usepackage{amsmath}
\usepackage{fontawesome5}
\usepackage{hyperref}
\hypersetup{
    colorlinks=true,
    linkcolor=blue,
    urlcolor=blue,
    citecolor=blue
}

% =====================================================================
% TEMA OG FARGER
% =====================================================================
\usetheme{Madrid}
\usecolortheme{beaver}

% =====================================================================
% TITTELINFO
% =====================================================================
\title{Neurosymbolic AI and Agentic AI}
\subtitle{BMED365: Topic List S01--S11}
\author{BMED365}
\date{Spring 2026}

% =====================================================================
% DOKUMENT
% =====================================================================
\begin{document}

% Title page
\begin{frame}
    \titlepage
\end{frame}

% Table of Contents
\begin{frame}{Overview}
    \begin{enumerate}
        \item \textbf{Fundamental paradigmer}
        \begin{itemize}
            \item S01: Kontrastere symbolsk og konneksjonistisk AI
            \item S02: Explain konseptet neurosymbolic integrasjon
        \end{itemize}
        \item \textbf{Knowledge Grapher og ontologier}
        \begin{itemize}
            \item S03: Describe hva en knowledge graph er
            \item S04: Know about medisinske ontologier
            \item S05: Discuss fordeler med neurosymbolic AI i medisin
        \end{itemize}
        \item \textbf{Case: Neurosymbolic gliomdiagnose}
        \begin{itemize}
            \item S06: Neurosymbolic AI for gliomdiagnostikk
            \item S07: Explain hvordan knowledge grapher kan validere prediksjoner
        \end{itemize}
        \item \textbf{Agentic AI}
        \begin{itemize}
            \item S08: Define agentic AI og dens kjerneegenskaper
            \item S09: Describe hvordan en AI-agent kan orkestrere klinisk arbeidsflyt
            \item S10--S11: Human-in-the-loop og etiske utfordringer
        \end{itemize}
    \end{enumerate}
\end{frame}

% =====================================================================
% SEKSJON: Symbolsk vs. Konneksjonistisk AI
% =====================================================================
\section{Fundamental paradigmer}

\begin{frame}{S01: Kontrastere symbolsk og konneksjonistisk AI}
    \begin{columns}[T]
        \begin{column}{0.48\textwidth}
            \textbf{\href{https://en.wikipedia.org/wiki/Symbolic_artificial_intelligence}{Symbolsk AI} (GOFAI):}
            \begin{itemize}
                \item ``Good Old-Fashioned AI''
                \item Eksplisitte regler og symboler
                \item Logisk resonnering
                \item Kunnskapsrepresentasjon
                \item Ekspert-systemer
            \end{itemize}
            
            \vspace{0.3cm}
            \textbf{Styrker:}
            \begin{itemize}
                \item Forklarbar
                \item Presise resonnementer
                \item Kodifiserer ekspertkunnskap
            \end{itemize}
        \end{column}
        \begin{column}{0.48\textwidth}
            \textbf{\href{https://en.wikipedia.org/wiki/Connectionism}{Konneksjonistisk AI}:}
            \begin{itemize}
                \item Neural Networks, deep learning
                \item Lærer fra data
                \item Subsymbolsk representasjon
                \item Mønstergjenkjenning
                \item ``Black box''
            \end{itemize}
            
            \vspace{0.3cm}
            \textbf{Styrker:}
            \begin{itemize}
                \item Lærer komplekse mønstre
                \item Håndterer støy og usikkerhet
                \item Skalerer med data
            \end{itemize}
        \end{column}
    \end{columns}
    
    \vspace{0.15cm}
    \begin{block}{Historisk perspektiv}
        AI har svingt mellom disse paradigmene -- se \href{https://en.wikipedia.org/wiki/AI_winter}{AI-vintrene}. Nå: Kan vi kombinere dem?
    \end{block}
\end{frame}

\begin{frame}{S02: Explain konseptet neurosymbolic integrasjon}
    \textbf{\href{https://en.wikipedia.org/wiki/Neuro-symbolic_AI}{Neurosymbolic AI} = det beste fra to verdener}
    
    \vspace{0.2cm}
    \textbf{Hovedidé:}
    \begin{itemize}
        \item Kombiner \textbf{nevralt} (læring fra data, mønstergjenkjenning)
        \item med \textbf{symbolsk} (resonnering, kunnskap, explainability)
    \end{itemize}
    
    \vspace{0.2cm}
    \textbf{Integrasjonsmønstre:}
    \begin{enumerate}
        \item \textbf{Neural $\rightarrow$ Symbolic:} Nevralt nettverk produserer symboler for resonnering
        \item \textbf{Symbolic $\rightarrow$ Neural:} Symbolsk kunnskap guider nevralt nettverk
        \item \textbf{Hybrid:} Tett integrasjon der begge informerer hverandre
    \end{enumerate}
    
    \vspace{0.2cm}
    {\small
    \begin{block}{Example: Medical diagnose (hjernetumor)}
        \textbf{Nevral:} CNN analyserer MR-bilde, identifiserer tumor-features \\
        \textbf{Symbolsk:} Regelbasert klassifikasjon iht. \href{https://www.who.int/publications/i/item/9789240021310}{WHO-kriterier} \\
        \textbf{Integrert:} Kombinerer bildefunn med kliniske regler for diagnose
    \end{block}
    }
\end{frame}

% =====================================================================
% SEKSJON: Knowledge Grapher
% =====================================================================
\section{Knowledge Grapher og ontologier}

\begin{frame}{S03: Describe hva en knowledge graph er}
    \textbf{\href{https://en.wikipedia.org/wiki/Knowledge_graph}{Knowledge Graph} (Knowledge Graph):}
    \begin{itemize}
        \item Strukturert representasjon av kunnskap som en \textbf{graf}
        \item Noder = \textbf{entiteter} (personer, sykdommer, legemidler)
        \item Kanter = \textbf{relasjoner} (``behandles med'', ``forårsaker'')
    \end{itemize}
    
    \vspace{0.1cm}
    \textbf{Trippel-representasjon:}
    \begin{center}
        (Subjekt, Predikat, Objekt) -- \textit{(Diabetes, behandles\_med, Metformin)}
    \end{center}
    
    \vspace{0.1cm}
    \begin{columns}[T]
        \begin{column}{0.48\textwidth}
            \textbf{Applications:}
            \begin{itemize}
                \item \href{https://en.wikipedia.org/wiki/Google_Knowledge_Graph}{Google Knowledge Graph}
                \item Medical beslutningsstøtte
                \item Legemiddelinteraksjoner
            \end{itemize}
        \end{column}
        \begin{column}{0.48\textwidth}
            \textbf{Advantages:}
            \begin{itemize}
                \item Maskinlesbar kunnskap
                \item Resonnering over relasjoner
                \item \href{https://en.wikipedia.org/wiki/Linked_data}{Linked Data}
            \end{itemize}
        \end{column}
    \end{columns}
    
    \vspace{0.05cm}
    {\scriptsize
    \begin{block}{Teknologier}
        \href{https://en.wikipedia.org/wiki/Resource_Description_Framework}{RDF}, \href{https://en.wikipedia.org/wiki/Web_Ontology_Language}{OWL}, \href{https://en.wikipedia.org/wiki/SPARQL}{SPARQL}, \href{https://en.wikipedia.org/wiki/SHACL}{SHACL} -- standarder for kunnskapsrepresentasjon, spørring og validation
    \end{block}
    }
\end{frame}

\begin{frame}{S04: Know about medisinske ontologier}
    \textbf{\href{https://en.wikipedia.org/wiki/Ontology_(information_science)}{Ontologi} = formell spesifikasjon av begreper og relasjoner i et domene}
    
    \vspace{0.3cm}
    \textbf{Importante medisinske ontologier/terminologier:}
    \begin{center}
        \footnotesize
        \begin{tabular}{lp{5cm}l}
            \toprule
            \textbf{Navn} & \textbf{Describelse} & \textbf{Application} \\
            \midrule
            \href{https://www.snomed.org/}{SNOMED CT} & Clinicale termer, 350 000+ begreper & \href{https://en.wikipedia.org/wiki/Electronic_health_record}{EPJ}, classification \\
            \href{https://www.who.int/standards/classifications/classification-of-diseases}{ICD-10/11} & Internasjonal sykdomsklassifikasjon & Rapportering, statistikk \\
            \href{https://loinc.org/}{LOINC} & Laboratorieundersøkelser & Labresultater \\
            \href{https://www.nlm.nih.gov/research/umls/rxnorm/}{RxNorm} & Legemidler & Forskrivning \\
            \midrule
            \href{https://www.who.int/publications/i/item/9789240021310}{WHO CNS 2021} & CNS-tumorklassifikasjon & Gliomdiagnose \\
            \bottomrule
        \end{tabular}
    \end{center}
    
    \vspace{0.3cm}
    \begin{block}{\href{https://www.who.int/publications/i/item/9789240021310}{WHO 2021 CNS-klassifikasjon}}
        Integrerer histologiske og molekylære kriterier. Example: \href{https://en.wikipedia.org/wiki/IDH1}{IDH-mutant} \href{https://en.wikipedia.org/wiki/Astrocytoma}{astrocytom} krever både histologi \textbf{og} IDH-mutasjonsstatus.
    \end{block}
\end{frame}

\begin{frame}{S05: Discuss fordeler med neurosymbolic AI i medisin}
    \textbf{Potensielle fordeler:}
    
    \vspace{0.3cm}
    \begin{enumerate}
        \item \textbf{\href{https://en.wikipedia.org/wiki/Explainable_artificial_intelligence}{Explainability}:} Symbolsk komponent gir tolkbare resonnementer
        
        \item \textbf{Dataholdighet:} Symbolsk kunnskap kompenserer for lite data
        
        \item \textbf{\href{https://en.wikipedia.org/wiki/Robustness_(computer_science)}{Robustness}:} Regelbaserte begrensninger forhindrer absurde prediksjoner
        
        \item \textbf{Oppdaterbarhet:} Ny medisinsk kunnskap legges til som regler/fakta
        
        \item \textbf{Samsvar med retningslinjer:} Kode WHO-kriterier, guidelines direkte
        
        \item \textbf{Validation:} Sjekk om prediksjoner er konsistente med kjent kunnskap
    \end{enumerate}
    
    \vspace{0.3cm}
    \begin{alertblock}{Why dette er viktig i medisin}
        Medisin er et domene der ekspertkunnskap er rik, explainability er kritisk, og feil kan være fatale. Ren deep learning er ofte utilstrekkelig.
    \end{alertblock}
\end{frame}

% =====================================================================
% SEKSJON: Case: Gliomdiagnose
% =====================================================================
\section{Case: Neurosymbolic gliomdiagnose}

\begin{frame}{S06: Neurosymbolic AI for gliomdiagnostikk}
    \textbf{Case study fra Lab 3, Notebook 08:}
    
    \vspace{0.3cm}
    \textbf{Problemet:}
    \begin{itemize}
        \item Gliom-klassifikasjon krever både \textbf{image analysis} (MRI) og \textbf{molekylære markører} (IDH, MGMT, 1p/19q)
        \item WHO 2021 krever integrasjon av flere informasjonskilder
    \end{itemize}
    
    \vspace{0.3cm}
    \textbf{Neurosymbolic løsning:}
    \begin{enumerate}
        \item \textbf{Nevral komponent:} CNN segmenterer tumor i MRI (BraTS)
        \item \textbf{Symbolsk komponent:} WHO-kriterier kodet som knowledge graph
        \item \textbf{Integrasjon:} Regelbasert klassifikator bruker CNN-output + molekylære markører
    \end{enumerate}
    
    \vspace{0.3cm}
    \begin{block}{Resultat}
        Diagnose som er både \textbf{bildebasert} (CNN) og \textbf{kriteriekonform} (WHO) -- med forklarbar resonnering.
    \end{block}
\end{frame}

\begin{frame}{S07: Explain hvordan knowledge grapher kan validere nevro-onkologiske prediksjoner}
    \textbf{Validation av \href{https://en.wikipedia.org/wiki/Neuro-oncology}{nevro-onkologiske} prediksjoner med symbolsk kunnskap:}
    
    \vspace{0.2cm}
    \textbf{Scenario:}
    \begin{itemize}
        \item CNN predikerer: ``Glioblastom med IDH-mutasjon''
        \item Knowledge Graph inneholder: ``Glioblastom er per definisjon IDH-villtype (WHO 2021)''
        \item \textbf{Konflikt detektert!}
    \end{itemize}
    
    \vspace{0.15cm}
    \textbf{Validationsmekanismer:}
    \begin{enumerate}
        \item \textbf{Konsistenssjekk:} Er prediksjonen logisk konsistent?
        \item \textbf{Regelvalidation:} Oppfyller prediksjonen nødvendige kriterier?
        \item \textbf{Plausibilitetssjekk:} Er kombinasjonen av funn realistisk?
    \end{enumerate}
    
    {\small
    \begin{alertblock}{Verdi}
        Fanger opp potensielle CNN-feil \textbf{før} de når kliniker. Øker tilliten til systemet.
    \end{alertblock}
    }
\end{frame}

% =====================================================================
% SEKSJON: Agentic AI
% =====================================================================
\section{Agentic AI}

\begin{frame}{S08: Define agentic AI og dens kjerneegenskaper}
    \textbf{\href{https://en.wikipedia.org/wiki/Agentic_AI}{Agentic AI} = AI som handler autonomt mot mål}
    
    \vspace{0.2cm}
    \textbf{Kjerneegenskaper:}
    \begin{enumerate}
        \item \textbf{Autonomi:} Handler (gjerne iterativt) uten kontinuerlig menneskelig input
        \item \textbf{Målrettet:} Jobber mot definerte mål
        \item \textbf{Planlegging:} Bryter ned komplekse mål i deloppgaver
        \item \textbf{Verktøybruk:} Kan bruke eksterne verktøy (API, databaser)
        \item \textbf{Refleksjon:} Evaluerer egen fremgang, justerer strategi
        \item \textbf{Minne:} Husker kontekst over flere interaksjoner og iterasjoner
    \end{enumerate}
    
    \vspace{0.2cm}
    \textbf{LLM-baserte agenter:}
    \begin{itemize}
        \item LLM som ``hjerne'' og ``orkestrator'' -- resonnerer og planlegger
        \item Verktøy-kall (\href{https://platform.openai.com/docs/guides/function-calling}{function calling}) for handlinger
        \item Examples: \href{https://github.com/Significant-Gravitas/AutoGPT}{AutoGPT}, \href{https://www.langchain.com/}{LangChain}, \href{https://docs.anthropic.com/claude/docs/tool-use}{Claude tool use}, \href{https://openai.com/index/introducing-gpt-4o-with-computer-use/}{OpenAI Operator}, \href{https://www.crewai.com/}{CrewAI}
    \end{itemize}
\end{frame}

\begin{frame}{S09: Describe hvordan en AI-agent kan orkestrere klinisk arbeidsflyt}
    \textbf{Example: Agentic AI for gliomutredning}
    
    \vspace{0.2cm}
    \textbf{Scenario:} Lege ber om ``Utred ny MR for mulig gliom''
    
    \vspace{0.2cm}
    \textbf{Agenten kan:}
    \begin{enumerate}
        \item \faSearch~\textbf{Hent data:} Koble til \href{https://en.wikipedia.org/wiki/Picture_archiving_and_communication_system}{PACS}, laste ned MR-bilder
        \item \faBrain~\textbf{Analyser:} Kjør \href{https://www.synapse.org/brats}{BraTS-segmentation}, beregn tumorvolum
        \item \faBookMedical~\textbf{Litteratur:} Søk \href{https://pubmed.ncbi.nlm.nih.gov/}{PubMed} for relevante studier
        \item \faDatabase~\textbf{Biobank:} Sjekk om molekylær analyse er tilgjengelig
        \item \faFile*~\textbf{Klassifiser:} Bruk WHO-kriterier for tentativ diagnose
        \item \faClipboard~\textbf{Rapporter:} Generer strukturert rapport for \href{https://en.wikipedia.org/wiki/Multidisciplinary_team}{MDT-møte}
    \end{enumerate}
    
    \vspace{0.1cm}
    {\small
    \begin{block}{Orkestrering}
        Agenten \textbf{koordinerer} flere AI-systemer og datakilder mot ett klinisk mål -- uten at legen må gjøre hvert steg manuelt.
    \end{block}
    }
\end{frame}

\begin{frame}{S10: Explain konseptet human-in-the-loop i agentiske systemer}
    \textbf{\href{https://en.wikipedia.org/wiki/Human-in-the-loop}{Human-in-the-loop} (HITL) i agentic AI:}
    
    \vspace{0.3cm}
    \begin{itemize}
        \item Selv autonome agenter trenger \textbf{\href{https://en.wikipedia.org/wiki/Human_oversight_of_AI}{menneskelig tilsyn}}
        \item Spesielt viktig for beslutninger med store konsekvenser
    \end{itemize}
    
    \vspace{0.3cm}
    \textbf{Implementeringsmønstre:}
    \begin{enumerate}
        \item \textbf{Godkjenningspunkter:} Agent pauser for godkjenning før kritiske handlinger
        \item \textbf{Konfidensterskel:} Menneske involveres når agent er usikker
        \item \textbf{\href{https://en.wikipedia.org/wiki/Audit_trail}{Audit trail}:} Alle handlinger logges for review
        \item \textbf{\href{https://en.wikipedia.org/wiki/Override}{Overstyring} (Overriding):} Menneske kan avbryte or korrigere agent
    \end{enumerate}
    
    \vspace{0.3cm}
    \begin{alertblock}{In medicinesk kontekst}
        Agent kan samle data og foreslå diagnose, men \textbf{legen} tar endelig beslutning. Agenten er en ``superkraftig assistent'', ikke en autonom beslutningstaker.
    \end{alertblock}
\end{frame}

\begin{frame}{S11: Discuss etiske utfordringer med autonome AI-agenter i helsevesenet}
    \textbf{Etiske utfordringer:}
    
    \vspace{0.3cm}
    \begin{enumerate}
        \item \textbf{Responsibility:} Hvem er responsibilitylig når en agent gjør feil?
        \begin{itemize}
            \item Utvikler? Sykehus? Lege som ``slapp løs'' agenten?
        \end{itemize}
        
        \item \textbf{Autonomi og kontroll:} Hvor mye autonomi er for mye?
        \begin{itemize}
            \item Risiko for utilsiktet handling
        \end{itemize}
        
        \item \textbf{Transparens:} Kan vi understand hvorfor agenten tok beslutningene?
        
        \item \textbf{Sikkerhet:} Agenter med tilgang til systemer = angrepsflate
        
        \item \textbf{Dehumanisering:} Risiko for å redusere menneskelig kontakt i helse
        
        \item \textbf{Overtillit:} ``\href{https://en.wikipedia.org/wiki/Automation_bias}{Automation bias}'' -- blindt stole på agenten
    \end{enumerate}
    
    \vspace{0.3cm}
    \begin{block}{Prinsipper for responsibilitylig agentic AI}
        Gradvis autonomi, robust HITL, klar responsibilitysfordeling, omfattende testing, kontinuerlig overvåking
    \end{block}
\end{frame}

% =====================================================================
% OPPSUMMERING
% =====================================================================
\section*{Summary}

\begin{frame}{Summary: Neurosymbolic AI and Agentic AI}
    \textbf{Neurosymbolic AI:}
    \begin{itemize}
        \item S01--S02: Kombinerer nevralt (læring) og symbolsk (resonnering)
        \item S03--S04: Knowledge Grapher og medisinske ontologier
        \item S05--S07: Advantages i medisin, case gliomdiagnose, validation
    \end{itemize}
    
    \vspace{0.3cm}
    \textbf{Agentic AI:}
    \begin{itemize}
        \item S08: Autonome, målrettede AI-systemer med planlegging
        \item S09: Orkestrering av kliniske arbeidsflyter
        \item S10--S11: HITL og etiske utfordringer
    \end{itemize}
    
    \vspace{0.3cm}
    \begin{block}{Lab 3, Notebook 08}
        Utforsk en detaljert case study som demonstrerer neurosymbolic gliomklassifikasjon og agentisk orkestrering.
    \end{block}
\end{frame}

\end{document}







